\documentclass[11pt,a4paper,fontset=fandol]{ctexart}

\usepackage[a4paper,top=2.25cm,bottom=2.45cm,left=2.25cm,right=2.25cm,headheight=18pt,headsep=0.55cm,footskip=24pt]{geometry}
\usepackage{amsmath,amssymb,amsthm}
\usepackage{fancyhdr}
\usepackage{lastpage}
\usepackage{enumitem}
\usepackage{titlesec}
\usepackage{xcolor}
\usepackage{hyperref}
\usepackage{bookmark}
\usepackage[most]{tcolorbox}
\usepackage{setspace}
\usepackage{array}

\hypersetup{
  colorlinks=true,
  linkcolor=blue!50!black,
  urlcolor=blue!50!black
}

\setstretch{1.13}
\setlength{\parindent}{2em}
\setlength{\parskip}{0.25em}
\allowdisplaybreaks
\setlist[enumerate]{itemsep=0.45em, topsep=0.35em, leftmargin=2.4em}
\setlist[itemize]{itemsep=0.25em, topsep=0.25em, leftmargin=1.9em}

\definecolor{mainblue}{RGB}{36,83,140}
\definecolor{lightblue}{RGB}{241,246,252}
\definecolor{lightgray}{RGB}{246,246,246}
\definecolor{rulegray}{RGB}{110,118,128}

\titleformat{\section}
  {\Large\bfseries\color{mainblue}}
  {\thesection.}
  {0.45em}
  {}
  [\vspace{0.2em}\color{rulegray}\titlerule]
\titlespacing*{\section}{0pt}{1.1em}{0.7em}

\pagestyle{fancy}
\fancyhf{}
\fancyhead[L]{\small 函数图像平移周测 A 卷}
\fancyhead[C]{\small 代数专题：基础与标准变式}
\fancyhead[R]{\small 刘文博}
\fancyfoot[C]{\small 第 \thepage 页 / 共 \pageref{LastPage} 页}
\renewcommand{\headrulewidth}{0.55pt}
\renewcommand{\footrulewidth}{0.35pt}

\newtcolorbox{infobox}[1]{
  breakable,
  colback=lightblue,
  colframe=mainblue,
  boxrule=0.7pt,
  arc=1mm,
  title=\textbf{#1},
  fonttitle=\bfseries
}

\newenvironment{solution}{\par\noindent\textbf{参考解答：}}{\hfill$\square$\par}
\newcommand{\answerspace}{\vspace{2.3cm}}
\newcommand{\biganswerspace}{\vspace{3.1cm}}

\begin{document}

\thispagestyle{empty}
\begin{center}
{\fontsize{22}{28}\selectfont\bfseries 函数图像平移周测 A 卷\par}
\vspace{0.4cm}
{\large 适合初中阶段函数专题基础检测\par}
\end{center}

\begin{infobox}{考试说明}
\begin{itemize}
  \item 测试时间：60--70 分钟
  \item 满分：100 分
  \item A 卷侧重标准图像平移与反向判断，先分清“上下改外面，左右改里面”。
\end{itemize}
\end{infobox}

\section{试题}

\begin{enumerate}
  \item （10 分）把函数
  \[
  y=x^2
  \]
  向上平移 $4$ 个单位，再向右平移 $3$ 个单位，求新方程。
  \answerspace

  \item （12 分）把函数
  \[
  y=|x|
  \]
  向左平移 $2$ 个单位，再向下平移 $3$ 个单位，求新方程。
  \answerspace

  \item （12 分）把函数
  \[
  y=2x-1
  \]
  向右平移 $4$ 个单位，再向上平移 $5$ 个单位，求化简后的新方程。
  \answerspace

  \item （12 分）已知函数
  \[
  y=(x-5)^2+2,
  \]
  说出它是由
  \[
  y=x^2
  \]
  怎样平移得到的。
  \answerspace

  \item （12 分）已知函数
  \[
  y=x^2+4x+1,
  \]
  说出它是由
  \[
  y=x^2
  \]
  怎样平移得到的。
  \answerspace

  \item （14 分）已知函数
  \[
  y=f(x)
  \]
  经过平移后变成
  \[
  y=f(x-5)+3.
  \]
  说出图像的平移方向和距离。
  \biganswerspace

  \item （14 分）已知抛物线
  \[
  y=x^2
  \]
  经过平移后，顶点变为 $(2,-1)$。求新方程。
  \biganswerspace

  \item （14 分）把函数
  \[
  y=|x-1|+2
  \]
  向左平移 $3$ 个单位，再向下平移 $4$ 个单位，求新方程，并写出新图像的顶点坐标。
  \biganswerspace
\end{enumerate}

\newpage
\section{详细解答}

\begin{enumerate}
  \item \begin{solution}
  先向上平移 $4$ 个单位，得到
  \[
  y=x^2+4.
  \]
  再向右平移 $3$ 个单位，应把 $x$ 换成 $x-3$，得到
  \[
  y=(x-3)^2+4.
  \]
  所以新方程为
  \[
  \boxed{y=(x-3)^2+4}.
  \]
  \end{solution}

  \item \begin{solution}
  向左平移 $2$ 个单位，应把 $x$ 换成 $x+2$，得
  \[
  y=|x+2|.
  \]
  再向下平移 $3$ 个单位，得
  \[
  \boxed{y=|x+2|-3}.
  \]
  \end{solution}

  \item \begin{solution}
  先向右平移 $4$ 个单位：
  \[
  y=2(x-4)-1.
  \]
  再向上平移 $5$ 个单位：
  \[
  y=2(x-4)-1+5.
  \]
  化简得
  \[
  y=2x-4.
  \]
  所以新方程为
  \[
  \boxed{y=2x-4}.
  \]
  \end{solution}

  \item \begin{solution}
  在
  \[
  y=(x-5)^2+2
  \]
  中，$x-5$ 表示向右平移 $5$ 个单位，外面的 $+2$ 表示向上平移 $2$ 个单位。

  因此它是由 $y=x^2$ \textbf{向右平移 $5$ 个单位，再向上平移 $2$ 个单位}得到的。
  \end{solution}

  \item \begin{solution}
  先配方：
  \[
  y=x^2+4x+1=(x+2)^2-3.
  \]
  所以它是由 $y=x^2$ \textbf{向左平移 $2$ 个单位，再向下平移 $3$ 个单位}得到的。
  \end{solution}

  \item \begin{solution}
  方程
  \[
  y=f(x-5)+3
  \]
  已经写成统一形式
  \[
  y=f(x-a)+b.
  \]
  这里 $a=5$，$b=3$。

  所以图像是由原图像\textbf{向右平移 $5$ 个单位，再向上平移 $3$ 个单位}得到的。
  \end{solution}

  \item \begin{solution}
  函数 $y=x^2$ 的顶点是 $(0,0)$。

  现在顶点变成 $(2,-1)$，说明图像向右平移了 $2$ 个单位，再向下平移了 $1$ 个单位。

  因此新方程为
  \[
  \boxed{y=(x-2)^2-1}.
  \]
  \end{solution}

  \item \begin{solution}
  先向左平移 $3$ 个单位，应把 $x$ 换成 $x+3$：
  \[
  y=|(x+3)-1|+2=|x+2|+2.
  \]
  再向下平移 $4$ 个单位：
  \[
  y=|x+2|+2-4=|x+2|-2.
  \]
  所以新方程为
  \[
  \boxed{y=|x+2|-2}.
  \]

  原函数 $y=|x-1|+2$ 的顶点是 $(1,2)$。

  向左平移 $3$ 个单位后变成 $(-2,2)$；再向下平移 $4$ 个单位后变成
  \[
  \boxed{(-2,-2)}.
  \]
  \end{solution}
\end{enumerate}

\end{document}
